\section{World-Frame Projection and Position Integration}
\label{sec:projection}

% ----------------------------------------------------------------
% TARGET: ~1 page. From quaternion to position.
% ----------------------------------------------------------------

\subsection{Rotating Acceleration to the World Frame}
\label{sec:proj:rotation}

Once the attitude quaternion $\mathbf{q}(t)$ has been estimated by the
complementary filter, each raw accelerometer reading
$\mathbf{a}_b(t)$ can be transformed into the world frame:

\begin{equation}
  \mathbf{a}_w(t) = \mathbf{q}(t) \otimes \mathbf{a}_{b,q}(t)
                    \otimes \mathbf{q}^*(t)
  \label{eq:a_world}
\end{equation}

\noindent where $\mathbf{a}_{b,q} = (0, a_x, a_y, a_z)$ is the accelerometer
reading embedded as a pure quaternion. The vector part of the result gives
$\mathbf{a}_w = (a_{w,x},\; a_{w,y},\; a_{w,z})$ in the world frame.

% ----- FIGURE: rotation illustration -----
% TODO: Illustration showing:
%   - A tilted sensor with a_body vector
%   - The quaternion rotation arrow
%   - The resulting a_world vector in the world frame
%   - Gravity subtraction step
%
% \begin{figure}[t]
%   \centering
%   \includegraphics[width=\columnwidth]{frame_rotation.pdf}
%   \caption{Projection of the body-frame acceleration into the world frame
%            via quaternion conjugation, followed by gravity removal.}
%   \label{fig:frame_rotation}
% \end{figure}

\subsection{Gravity Removal}
\label{sec:proj:gravity}

In the world frame, gravity appears as a constant vector
$\mathbf{g}_w = (0,\; 0,\; g)^\top$ (assuming ENU convention with
$g \approx 9.81\;\text{m/s}^2$). The linear acceleration is:

\begin{equation}
  \mathbf{a}_\text{lin}(t) = \mathbf{a}_w(t) - \mathbf{g}_w
  \label{eq:gravity_removal}
\end{equation}

\noindent This step is possible \emph{only} because the rotation to the world
frame has been performed first. In the body frame, the gravity component
changes direction with every rotation, making subtraction impossible.

\subsection{Double Integration to Position}
\label{sec:proj:integration}

Velocity and position are obtained by numerical integration of the
gravity-free acceleration:

\begin{align}
  \mathbf{v}(t + \Delta t) &= \mathbf{v}(t)
    + \mathbf{a}_\text{lin}(t)\,\Delta t
  \label{eq:velocity} \\[4pt]
  \mathbf{p}(t + \Delta t) &= \mathbf{p}(t)
    + \mathbf{v}(t)\,\Delta t
    + \tfrac{1}{2}\,\mathbf{a}_\text{lin}(t)\,\Delta t^2
  \label{eq:position}
\end{align}

\noindent with initial conditions $\mathbf{v}(0) = \mathbf{0}$ and
$\mathbf{p}(0) = \mathbf{0}$ (the sensor starts at rest at the origin).

\subsection{Drift in Position Estimation}
\label{sec:proj:drift}

Double integration amplifies any residual error in
$\mathbf{a}_\text{lin}$: a constant bias $\epsilon$ in acceleration produces
a position error that grows as $\frac{1}{2}\epsilon\,t^2$. This error has two
distinct sources:

\paragraph{Static bias.}
The accelerometer has a small, approximately constant offset that persists
even at rest. This can be estimated and removed during an initial calibration
phase by averaging the linear acceleration over a stationary period of
several seconds.

\paragraph{Gravity leakage from orientation error.}
This is the dominant source of drift. Even a small error $\delta$ in the
estimated roll or pitch causes the gravity vector to be imperfectly
subtracted after the world-frame rotation, leaving a residual:

\begin{equation}
  \epsilon_g \approx g \sin(\delta)
  \label{eq:gravity_leakage}
\end{equation}

\noindent For example, an orientation error of just $\delta = 1.4°$ produces
a residual of $9.81 \times \sin(1.4°) \approx 0.24$~m/s$^2$. After double
integration over a time interval $T$:

\begin{equation}
  \Delta p = \tfrac{1}{2}\,\epsilon_g\,T^2
           = \tfrac{1}{2}\,g\sin(\delta)\,T^2
  \label{eq:drift_position}
\end{equation}

\noindent For $T = 3$~s, this yields $\Delta p \approx 1.08$~m of spurious
displacement. Crucially, this error \emph{changes direction} as the device
rotates---it is not a fixed offset that can be calibrated away. It affects all
three world-frame axes (X, Y, Z), not just the vertical. This is the
fundamental limitation of inertial navigation with consumer-grade MEMS
sensors.

\subsection{Drift Mitigation Strategies}
\label{sec:proj:mitigation}

Three complementary strategies are applied to reduce the drift:

\subsubsection{Static Calibration with Measured Gravity}
\label{sec:proj:calibration}

Rather than using the textbook value $g = 9.81$~m/s$^2$, the actual gravity
magnitude is estimated from the sensor data during the initial rest phase:

\begin{equation}
  g_\text{meas} = \frac{1}{N_c} \sum_{n=1}^{N_c}
    \lVert \mathbf{a}_b[n] \rVert
  \label{eq:g_measured}
\end{equation}

\noindent where $N_c$ is the number of samples in the rest phase.
The residual linear acceleration bias is then estimated as the mean of
$\mathbf{a}_\text{lin}$ over the same interval and subtracted from the entire
signal. Using the full rest phase ($N_c \approx 1000$ samples at 100~Hz,
corresponding to $\sim$10~s) instead of a handful of samples dramatically
improves the bias estimate.

\subsubsection{Zero-Velocity Updates (ZUPT)}
\label{sec:proj:zupt}

When the sensor is stationary, its true velocity is zero. Any non-zero
integrated velocity at that instant is, by definition, accumulated error.
A \emph{zero-velocity update} exploits this constraint by resetting the
velocity to zero whenever the sensor is detected to be at rest.

\paragraph{Stationary detection.}
Two conditions must hold simultaneously over a sliding window of $W$
consecutive samples:

\begin{align}
  \bigl|\,\lVert \mathbf{a}_b[n] \rVert - g\,\bigr|
    &< \sigma_a
  \label{eq:zupt_accel} \\[4pt]
  \lVert \boldsymbol{\omega}_b[n] \rVert
    &< \sigma_\omega
  \label{eq:zupt_gyro}
\end{align}

\noindent The first condition checks that the accelerometer norm is close to
gravity (no linear acceleration). The second checks that the angular velocity
is near zero (no rotation). Requiring both conditions over a window of $W$
samples (typically $W = 5$) prevents false positives from momentary pauses in
a single channel. Typical thresholds are $\sigma_a = 0.3$~m/s$^2$ and
$\sigma_\omega = 0.15$~rad/s.

\paragraph{Velocity reset.}
At every sample classified as stationary, the integrated velocity is forced to
zero on all three axes:

\begin{equation}
  \text{if stationary}[n]: \quad \mathbf{v}[n] = \mathbf{0}
  \label{eq:zupt_reset}
\end{equation}

\noindent This prevents unbounded velocity drift during rest periods and
provides anchor points for the subsequent backward correction.

\subsubsection{Backward Velocity Correction}
\label{sec:proj:backward}

ZUPT alone resets velocity at rest, but during movement segments the velocity
still drifts due to gravity leakage. The \emph{backward correction} addresses
this by distributing the accumulated velocity error linearly across each
movement segment.

\paragraph{Segment identification.}
A movement segment is defined as the interval between two consecutive
stationary phases. Let $n_s$ be the first moving sample (where the sensor
leaves rest) and $n_e$ be the last moving sample (just before the next rest
phase begins). At $n_e$, the velocity $\mathbf{v}[n_e]$ should ideally
transition smoothly to zero, but in practice it carries an error
$\mathbf{e}_v = \mathbf{v}[n_e]$ due to drift accumulated during the segment.

\paragraph{Linear error distribution.}
The error is assumed to grow approximately linearly during the segment (since
gravity leakage acts as a roughly constant bias over short intervals). It is
therefore distributed proportionally:

\begin{equation}
  \mathbf{v}_\text{corr}[n] = \mathbf{v}[n]
    - \mathbf{e}_v \cdot \frac{n - n_s + 1}{n_e - n_s + 1},
    \quad n \in [n_s,\, n_e]
  \label{eq:backward_correction}
\end{equation}

\noindent At $n = n_s$ the correction is minimal; at $n = n_e$ the full error
is subtracted, ensuring a smooth transition to zero velocity at the next ZUPT.
The corrected velocity is then integrated to obtain the position estimate.

\paragraph{Limitations.}
The backward correction is effective when the user's motion naturally includes
pauses (e.g.\ drawing shapes with intermittent stops). However, during
continuous uninterrupted motion, no ZUPT events occur and the correction
cannot be applied. This motivates the high-pass filtering strategy described
next.

\subsubsection{High-Pass Filtering of Linear Acceleration}
\label{sec:proj:highpass}

Even after gravity subtraction and static bias removal, the linear
acceleration retains a slowly varying residual caused by orientation
estimation errors. As shown in
Section~\ref{sec:proj:drift}, an orientation error of $\delta$ produces a
gravity leakage $\epsilon_g \approx g\sin(\delta)$ that changes direction as
the device rotates but maintains a \emph{non-zero mean} over time. This
non-zero mean, when integrated twice, produces the characteristic
monotonic drift observed in position estimates.

The key insight is that \textbf{real translational acceleration has zero
mean}: every acceleration phase (push) is followed by a deceleration phase
(stop), so the integral over a complete movement cycle is approximately zero.
In contrast, gravity leakage is a low-frequency bias with non-zero mean. A
high-pass filter can therefore separate the two.

\paragraph{First-order IIR high-pass filter.}
We apply a causal first-order high-pass filter independently to each axis of
$\mathbf{a}_\text{lin}$:

\begin{equation}
  y[n] = \alpha_\text{hp}\,\bigl(y[n-1] + x[n] - x[n-1]\bigr)
  \label{eq:highpass}
\end{equation}

\noindent where $x[n]$ is the input (linear acceleration after bias removal)
and $y[n]$ is the filtered output. The parameter
$\alpha_\text{hp} \in (0,\,1)$ controls the cutoff frequency:

\begin{equation}
  f_c = \frac{1 - \alpha_\text{hp}}{2\pi\,\Delta t}
  \label{eq:hp_cutoff}
\end{equation}

\noindent Signals below $f_c$ are attenuated (including the gravity leakage
bias); signals above $f_c$ pass through (the actual motion).

\paragraph{Choosing the cutoff frequency.}
The cutoff must satisfy two constraints simultaneously:

\begin{enumerate}
  \item $f_c$ must be \textbf{above} the frequency of the gravity leakage
        bias, which is quasi-static ($\ll 1$~Hz).
  \item $f_c$ must be \textbf{below} the fundamental frequency of the actual
        motion, so that the real acceleration signal is preserved.
\end{enumerate}

\noindent The optimal value therefore depends on the \emph{type of motion},
not on the individual recording:

\begin{itemize}
  \item \textbf{Fast hand gestures} ($\sim$2--4~Hz): $\alpha_\text{hp}
        \approx 0.85$, giving $f_c \approx 2.4$~Hz at 100~Hz sampling.
  \item \textbf{Walking or slow movements} ($\sim$1--2~Hz):
        $\alpha_\text{hp} \approx 0.95$, giving $f_c \approx 0.8$~Hz.
  \item \textbf{Very slow gestures} ($< 1$~Hz): $\alpha_\text{hp}
        \approx 0.98$, giving $f_c \approx 0.3$~Hz. At this point the filter
        provides minimal benefit, as the leakage frequency overlaps with the
        motion frequency.
\end{itemize}

\noindent If $f_c$ is set too high, the filter removes part of the real
acceleration signal, causing the estimated displacement to be
\emph{underestimated}. If $f_c$ is set too low, the gravity leakage passes
through and drift reappears. In practice, $\alpha_\text{hp}$ does not need to
be re-tuned for each recording---it only needs to match the general category
of motion (fast vs.\ slow).

\paragraph{Effect on drift.}
In our experiments, adding the high-pass filter reduced the position drift
from several meters to the order of $\pm 0.2$~m for a 17-second recording of
repetitive up-and-down hand motion, even without any ZUPT corrections (no
pauses in the motion). The mean residual acceleration on the Z-axis dropped
from $+0.41$~m/s$^2$ (without filter) to $+0.02$~m/s$^2$ (with
$\alpha_\text{hp} = 0.85$).

\subsubsection{Summary of the Drift Mitigation Pipeline}
\label{sec:proj:pipeline_summary}

The complete pipeline applies five corrections in sequence:

\begin{enumerate}
  \item \textbf{Measured gravity magnitude}: $g_\text{meas}$ from the rest
        phase replaces the nominal $9.81$~m/s$^2$.
  \item \textbf{Static bias removal}: mean of
        $\mathbf{a}_\text{lin}$ during rest is subtracted.
  \item \textbf{Moving average}: a short window ($\sim$5 samples) smooths
        high-frequency noise and vibration spikes.
  \item \textbf{High-pass filter}: removes the low-frequency gravity leakage
        bias that causes monotonic drift.
  \item \textbf{Acceleration clamp}: limits $|\mathbf{a}_\text{lin}|$ to a
        physically reasonable maximum (e.g.\ $\pm 4$~m/s$^2$ for hand motion),
        rejecting residual spikes from large orientation errors.
\end{enumerate}

\noindent After filtering, ZUPT and backward velocity correction are applied
during integration when stationary phases are available. For continuous motion
without pauses, the high-pass filter alone provides the primary drift
mitigation.

% ----- FIGURE: full pipeline -----
% TODO: End-to-end block diagram:
%   [Accel] ──→ [Tilt] ──→ [LPF] ──┐
%   [Gyro]  ──→ [∫ω]  ──→ [HPF] ──┤→ [q(t)] → [Rotate a_b] → [-g] → [∫∫] → [p(t)]
%
% \begin{figure*}[t]
%   \centering
%   \includegraphics[width=\textwidth]{full_pipeline.pdf}
%   \caption{Complete signal-processing pipeline: from raw sensor data to
%            three-dimensional position. The complementary filter estimates
%            orientation; the quaternion rotation projects acceleration into
%            the world frame; double integration yields position.}
%   \label{fig:pipeline}
% \end{figure*}
